% 电子版入口。与 CUMCMThesis 示例手册分开维护。
\documentclass[withoutpreface]{cumcmthesis}
\title{待填写：论文题目}
\begin{document}
\maketitle
\begin{abstract}
待填写：正文和结果核验完成后反写摘要。
\keywords{待填写}
\end{abstract}

\section{问题重述}
待填写。
\section{问题分析}
待填写。
\section{模型假设}
待填写：仅采用团队确认的必要假设，说明依据与后文作用。
\section{符号说明}
% 从实际公式提取全局主要符号，局部符号就近定义。
\begin{table}[htbp]
\centering
\caption{主要符号说明}
\label{tab:symbols}
\begin{tabularx}{0.95\textwidth}{cXc}
\toprule
符号 & 含义 & 单位 \\
\midrule
待填写 & 待填写 & 待填写 \\
\bottomrule
\end{tabularx}
\end{table}
\section{模型建立与求解}
待填写：每个结论关联真实结果版本。
\section{检验与灵敏度分析}
待填写。
\section{模型评价}
待填写。

% 必须在参考文献之前；声明与支撑材料由同一份台账生成。
% 由 scripts/ai-disclosure.py build 依据实际台账生成；不要手工编造已使用/未使用状态。
\section*{AI工具使用声明}
\textbf{待填写：请依据实际 AI 使用记录生成声明；当前为模板，不可提交。}

\begin{thebibliography}{99}
\bibitem{pending} 待填写：核实作者、题名、年份和 DOI / 来源后列入。
\end{thebibliography}

% 正文页数计至参考文献结束，包含AI声明；附录从新页开始。
\clearpage
\begin{appendices}
\section{支撑材料文件列表}
待填写：使用 AI 时列出 AI工具使用详情.pdf，归入支撑材料。
\section{软件环境与运行说明}
待填写：实际语言、库与求解器版本、启动目录、输入输出及复现命令。
\section{完整源程序}
待填写：按问题列出完整源码、公共模块和实际软件交互命令，与支撑包保持同版。
% 填入真实源码路径后可用 \lstinputlisting 引入，避免手工复制漂移。
\section{补充结果}
待填写：必要长表和补充检验；核心结果仍保留在正文。
\end{appendices}
\end{document}
